% capitulos/06-conclusao.tex

Lorem ipsum dolor sit amet, consectetur adipiscing elit. Os resultados
obtidos nesta pesquisa permitem concluir que a distribuição de mercúrio
total nos solos do Cerrado é controlada primariamente pelo teor de
matéria orgânica e apresenta variação sazonal significativa, com
concentrações mais elevadas no período de transição seca-chuva
\cite{exemplo:artigo:2024}.

A hipótese inicial foi corroborada pelos dados: as concentrações de HgT
variam significativamente com a profundidade do solo e com a proximidade
a fontes antrópicas de emissão \cite{exemplo:tese:2023}. O fator de
bioacumulação calculado para os organismos da fauna edáfica revelou
espécies sentinela potencialmente úteis para programas de biomonitoramento
\cite{exemplo:anais:2023}.

\section{Contribuições Originais}

As principais contribuições originais desta pesquisa são:

\begin{enumerate}
  \item Caracterização inédita da distribuição sazonal de HgT em
        três horizontes de solo em área de savana protegida do
        Distrito Federal;
  \item Identificação de correlação significativa entre HgT e COT
        em solos de Cerrado sob diferentes fitofisionomias;
  \item Proposição de índice de risco ecológico adaptado às
        condições pedológicas do Bioma Cerrado.
\end{enumerate}

\section{Limitações e Perspectivas Futuras}

Excepteur sint occaecat cupidatat non proident, sunt in culpa qui
officia deserunt mollit anim id est laborum \cite{exemplo:livro:2020}.
Entre as limitações do presente estudo destacam-se a ausência de dados
históricos anteriores a 2020 para a área de estudo e a não determinação
de metilmercúrio (MeHg), forma mais tóxica e biodisponível do elemento.

Como perspectivas futuras, recomenda-se a ampliação da rede de
monitoramento para outras fitofisionomias do Cerrado e a aplicação de
modelos de aprendizado de máquina para predição espacial das concentrações
de HgT, conforme proposto por \citeonline{exemplo:dissertacao:2022}.
